\documentclass[11pt,a4paper]{article}
\usepackage[margin=2.2cm]{geometry}
\usepackage{fontspec}
\usepackage{booktabs}
\usepackage{tabularx}
\usepackage{longtable}
\usepackage{hyperref}
\usepackage{xcolor}
\usepackage{enumitem}
\usepackage{titlesec}
\usepackage{amssymb}
\setmainfont{Times New Roman}
\definecolor{headingblack}{HTML}{000000}
\hypersetup{colorlinks=true, linkcolor=headingblack, urlcolor=headingblack}
\titleformat{\section}{\Large\bfseries\color{headingblack}}{\thesection}{0.8em}{}
\newcommand{\placeholder}[1]{[#1]}

\begin{document}

\begin{center}
{\LARGE\bfseries Supplementary Materials Template}\\[0.5em]
{\large Primary language: English. Use this document when extra files, code, data, links, screenshots, or videos are submitted.}\\[1.5em]
\end{center}

\begin{tabularx}{\textwidth}{>{\bfseries}p{0.24\textwidth}X}
\toprule
Project Title & \placeholder{Full project title}\\
Team Name & \placeholder{Team name}\\
Institution & \placeholder{Institution / school / organization}\\
Team Members & \placeholder{Names and roles}\\
Track / Category & \placeholder{If applicable}\\
Date & \placeholder{YYYY-MM-DD}\\
\bottomrule
\end{tabularx}

\section{Package Overview}
\placeholder{Briefly describe what is included in the supplementary package and how it supports the project report and technical abstract.}

\section{Recommended Directory Structure}
\begin{verbatim}
submission_package/
  report.pdf
  technical_abstract.pdf
  supplementary_materials.pdf
  code/
  data/
  results/
  figures/
  demo/
\end{verbatim}

\section{Code and Runtime Environment}
\begin{itemize}
  \item Operating system: \placeholder{Windows / Linux / macOS / container}
  \item Language and version: \placeholder{Python / R / JavaScript / etc.}
  \item Dependencies: \placeholder{requirements.txt, environment.yml, package.json, or Dockerfile}
  \item Entry point: \placeholder{Main script, notebook, or command}
  \item Expected runtime and hardware: \placeholder{CPU / GPU / memory}
\end{itemize}

\section{Data Description}
\begin{itemize}
  \item Data source and license.
  \item File format, sample count, feature count, and key fields.
  \item Preprocessing and quality control.
  \item Access restrictions or privacy notes.
\end{itemize}

\section{Results and Demonstration Materials}
\begin{itemize}
  \item Result files and how to interpret them.
  \item Screenshots, interface pages, or workflow captures.
  \item Platform URL, repository URL, demo video URL, and access credentials if applicable.
  \item Consistency note linking each result back to the report section.
\end{itemize}

\section{Reproducibility Checklist}
\begin{tabularx}{\textwidth}{p{0.08\textwidth}X X}
\toprule
\textbf{Check} & \textbf{Item} & \textbf{Notes}\\
\midrule
$\square$ & README describes the package structure. & \placeholder{Notes}\\
$\square$ & Main commands are provided and tested. & \placeholder{Notes}\\
$\square$ & Input and output paths are documented. & \placeholder{Notes}\\
$\square$ & Main parameters are listed. & \placeholder{Notes}\\
$\square$ & External data, model, or software licenses are stated. & \placeholder{Notes}\\
$\square$ & Supplementary claims match the report and abstract. & \placeholder{Notes}\\
\bottomrule
\end{tabularx}

\end{document}
